\documentclass[11pt]{article}
\usepackage[margin=1in]{geometry}
\usepackage{amsmath,amssymb,amsthm,mathtools,bm,booktabs,array,longtable}
\usepackage[hidelinks]{hyperref}
\usepackage{xcolor}
\usepackage{microtype}

\hypersetup{
  pdftitle={Fibered Renormalization as Representational Dynamics},
  pdfauthor={Research dossier},
  pdfsubject={Restricted Yang--Mills block-RG theorems and frontier audit},
  pdfkeywords={Yang--Mills, mass gap, renormalization group, heat kernel, conditional fiber}
}

\definecolor{openred}{RGB}{150,25,25}
\definecolor{provedgreen}{RGB}{20,105,55}
\definecolor{conditionalblue}{RGB}{25,65,145}
\newcommand{\proved}{\textcolor{provedgreen}{\textsc{proved}}}
\newcommand{\conditional}{\textcolor{conditionalblue}{\textsc{conditional}}}
\newcommand{\open}{\textcolor{openred}{\textsc{open}}}
\newcommand{\TV}{\operatorname{TV}}
\newcommand{\Cov}{\operatorname{Cov}}
\newcommand{\E}{\mathbb E}
\newcommand{\R}{\mathbb R}
\newcommand{\T}{\mathbb T}
\newcommand{\ind}{\mathbf 1}

\newtheorem{definition}{Definition}[section]
\newtheorem{theorem}[definition]{Theorem}
\newtheorem{proposition}[definition]{Proposition}
\newtheorem{lemma}[definition]{Lemma}
\newtheorem{corollary}[definition]{Corollary}
\theoremstyle{remark}
\newtheorem{remark}[definition]{Remark}

\title{Fibered Renormalization as Representational Dynamics:\\
An Interacting-Block Attack on the Yang--Mills Mass-Gap Programme}
\author{Research dossier---restricted theorems, calibrations, and frontier audit}
\date{20 July 2026}

\begin{document}
\maketitle

\begin{abstract}
We refine a representation-aware constructive programme for four-dimensional
pure Yang--Mills theory.  The central correction is that a renormalization step
must not be required to contract the degrees of freedom it intentionally
retains.  It should transport the retained coarse information and contract only
the conditional fiber that it eliminates.  A two-plaquette heat-kernel block
makes this distinction exact: in the Gaussian tangent model the uncentered
internal coordinate has boundary response $1/2$, whereas the centered fiber is
independent of the retained total holonomy and has response zero.  We prove an
exact multiscale covariance decomposition, formulate a conditional fibered
clustering theorem, prove an independent-defect connected-cluster bound, and
give a finite-product regulator comparison.  Compact $U(1)$ computations
calibrate these claims.  None of these results proves four-dimensional basin
entry, continuum existence, or a mass gap; those obligations are isolated
explicitly.
\end{abstract}

\begin{center}
\fbox{\parbox{0.92\textwidth}{\textbf{Status.}
The compact bridge and Gaussian statements, the covariance identity, the
independent-defect bound, and product-TV comparison are \proved.
The clustering implication is \conditional\ on stated estimates.
The Clay target remains \open.}}
\end{center}

\tableofcontents

\section{The target and the representational correction}

The Clay problem asks, for every compact simple gauge group $G$, for a
nontrivial quantum Yang--Mills theory on $\R^4$ satisfying axiomatic properties
and possessing a mass gap $\Delta>0$ \cite{jaffe-witten}.  In Euclidean language
a constructive route must produce coherent gauge-invariant Schwinger functions,
verify reflection positivity and the remaining reconstruction hypotheses, and
obtain decay strong enough to imply a positive spectral threshold after
Osterwalder--Schrader reconstruction \cite{osterwalder-schrader1,osterwalder-schrader2}.

The guiding ontology is not that one coordinate system is the theory.  Gauge
potentials are local coordinates on relational information; a regulator is a
temporary realization; a particle is a spectral interpretation after
reconstruction; renormalization changes descriptive scale; and the continuum
theory is coherence across compatible limiting representations.  On this view
the mass gap is an invariant form of infrared spectral rigidity.

The earlier one-step question ``does every block variable contract?'' is too
coarse.  An RG transformation separates a block into
\[
   x \longmapsto (y,r), \qquad y=\pi(x),
\]
where $y$ is retained and $r$ is integrated out.  Relevant information in $y$
need not contract; indeed its survival is the point of blocking.  A defensible
local question is instead:
\begin{quote}
After fixing the retained gauge-invariant coarse data and quotienting gauge
redundancy, does the conditional law on the eliminated fiber forget distant
boundary perturbations, apart from a quantitatively controlled defect channel?
\end{quote}

\section{A two-plaquette heat-kernel block}

Let $\T=\R/2\pi\mathbb Z$ and let
\[
 p_t(\theta)=\frac{1}{\sqrt{2\pi t}}
 \sum_{n\in\mathbb Z}\exp\!\left[-\frac{(\theta+2\pi n)^2}{2t}\right]
 =\frac1{2\pi}\sum_{k\in\mathbb Z}e^{-k^2t/2}e^{ik\theta}
\]
be the $U(1)$ heat kernel in the convention used here.  It satisfies the exact
convolution semigroup identity $p_t*p_u=p_{t+u}$.

Consider two adjacent plaquette angles $q_1,q_2$ with density
$p_t(q_1)p_t(q_2)$ and retain their total holonomy
\[
 s=q_1+q_2\pmod{2\pi}.
\]
The internal coordinate is $q=q_1$ and a centered fiber coordinate is
$r=q-s/2$ (locally, with the usual circle ambiguity away from the small-field
chart).

\begin{proposition}[Compact heat-kernel bridge, \proved]
For every $t>0$ and $s\in\T$, a regular conditional density of $q$ given $s$
is
\begin{equation}\label{eq:bridge}
 \gamma_t^s(q)=\frac{p_t(q)p_t(s-q)}{p_{2t}(s)}.
\end{equation}
Consequently the centered fiber has density
$\widetilde\gamma_t^s(r)=\gamma_t^s(r+s/2)$.
\end{proposition}
\begin{proof}
The joint density in coordinates $(q,s)$ is $p_t(q)p_t(s-q)$, because Haar
measure and the linear change of circle coordinates have unit Jacobian.  Its
$q$-integral is $(p_t*p_t)(s)=p_{2t}(s)$.  Division gives
\eqref{eq:bridge}; translation gives the centered formula.
\end{proof}

This is already an interacting conditional law: conditioning couples the two
plaquettes.  It is exactly solvable because the heat kernel is closed under
convolution, a feature familiar in two-dimensional Yang--Mills constructions
\cite{driver2d}.  It is a calibration model, not a substitute for the
four-dimensional nonabelian block.

\section{The retained-mode obstruction and its fiber resolution}

The small-field tangent model replaces the circle by $\R$.  Let
$X_1,X_2\stackrel{\mathrm{iid}}\sim N(0,t)$, retain $S=X_1+X_2$, and set
$R=X_1-S/2$.

\begin{theorem}[Exact Gaussian response split, \proved]\label{thm:gaussian}
For each $s\in\R$,
\[
 X_1\mid S=s\sim N(s/2,t/2),
 \qquad R\mid S=s\sim N(0,t/2).
\]
Therefore, for the $1$-Wasserstein metric,
\[
 \frac{W_1(\mathcal L(X_1\mid S=s),\mathcal L(X_1\mid S=s'))}{|s-s'|}
 =\frac12,
 \qquad
 \frac{W_1(\mathcal L(R\mid S=s),\mathcal L(R\mid S=s'))}{|s-s'|}=0.
\]
\end{theorem}
\begin{proof}
Gaussian conditioning gives mean
$\Cov(X_1,S)\operatorname{Var}(S)^{-1}s=(t/2t)s=s/2$ and conditional variance
$t-t^2/(2t)=t/2$.  Subtracting $s/2$ removes the conditional mean.  Translates
of the same one-dimensional probability measure have $W_1$ equal to the
translation distance, proving both ratios.
\end{proof}

The coefficient $1/2$ is not a failure of mixing: it is deterministic transport
of retained information.  Centering separates this transport from the
fluctuation that is actually eliminated.  On a longer Gaussian chain, analogous
interpolation of retained endpoint data can make a full-coordinate Dobrushin
row sum reach one even when the pinned fluctuation field is well controlled.
Thus full-mode contraction is not representation invariant.

For the compact bridge, winding sectors make
$\widetilde\gamma_t^s$ weakly $s$-dependent, but this dependence is
exponentially suppressed in a small-field chart as $t\downarrow0$.  Section
\ref{sec:numerics} verifies this numerically without promoting the observation
to a four-dimensional theorem.

\section{Fibered representation packages}

\begin{definition}[Fibered RG step]
A fibered RG step is a tuple
\[
 \mathfrak R_k=(X_k,\mu_k,X_{k+1},\pi_k,Q_k,\mathcal O_k,\mathcal G_k,\mathcal B_k)
\]
with:
\begin{enumerate}
\item a probability space $(X_k,\mu_k)$ and a coarse space $X_{k+1}$;
\item a measurable retained-data map $\pi_k:X_k\to X_{k+1}$ and
      $\mu_{k+1}=(\pi_k)_*\mu_k$;
\item a disintegration
      $\mu_k(dx)=\int Q_k(y,dx)\,\mu_{k+1}(dy)$ supported on
      $\pi_k^{-1}(y)$;
\item a gauge-invariant observable core $\mathcal O_k$ and an explicit
      transport rule $P_kF(y)=\int F(x)Q_k(y,dx)$;
\item a good event $\mathcal G_k$ on which conditional fiber influence is
      bounded by $K_k$, and a defect family $\mathcal B_k$ whose connected
      contribution is bounded by $R_k$.
\end{enumerate}
The retained base $y$ is transported.  The numbers $K_k$ and $R_k$ measure only
loss of memory through the conditional fiber and defects, respectively.
\end{definition}

For gauge theory, $X_k$ may be a regulated connection space modulo specified
local gauge actions; $y$ must consist of explicitly chosen coarse covariant or
invariant data.  This definition does not solve gauge fixing: existence,
measurability, locality, stabilizers, Gribov-type ambiguity, topology, and
reflection compatibility are part of the open construction of $\mathfrak R_k$.

\section{Exact covariance bookkeeping across scales}

Let $\mathcal F_{k+1}\subseteq\mathcal F_k$ be the sigma-algebra generated by
the retained map at step $k$, and write $P_kF=\E[F\mid\mathcal F_{k+1}]$.

\begin{theorem}[One-step fiber decomposition, \proved]\label{thm:one-step}
For $F,G\in L^2(\mu_k)$,
\begin{equation}\label{eq:total-cov}
 \Cov_{\mu_k}(F,G)
 =\Cov_{\mu_{k+1}}(P_kF,P_kG)
 +\E_{\mu_{k+1}}\!\left[
       \Cov_{Q_k(y,\cdot)}(F,G)
   \right].
\end{equation}
\end{theorem}
\begin{proof}
Insert $F=P_kF+(F-P_kF)$ and similarly for $G$.  The mixed terms vanish after
conditioning because $\E[F-P_kF\mid\mathcal F_{k+1}]=0$.  The covariance of
the conditional means is the first term; the conditional covariance of the
residuals is the second.
\end{proof}

\begin{corollary}[Iterated fiber decomposition, \proved]\label{cor:iterated}
For a tower $\mathcal F_0\supseteq\cdots\supseteq\mathcal F_n$, define
$F_k=\E[F\mid\mathcal F_k]$ and $G_k=\E[G\mid\mathcal F_k]$.  Then
\begin{equation}\label{eq:iterated}
 \Cov(F,G)=\Cov(F_n,G_n)
 +\sum_{k=0}^{n-1}
 \E\!\left[\Cov(F_k,G_k\mid\mathcal F_{k+1})\right].
\end{equation}
\end{corollary}
\begin{proof}
Apply Theorem~\ref{thm:one-step} to $(F_k,G_k)$ and telescope.
\end{proof}

Equation \eqref{eq:iterated} is the precise form of representational dynamics
used here.  Each step declares what information survives in $(F_{k+1},G_{k+1})$
and accounts separately for what is forgotten in the fiber covariance.  It is
an identity, not a decay estimate.

\section{A conditional fibered rigidity theorem}

We now state an abstract implication that makes the missing estimates visible.
Let $d_k(A,B)$ denote separation of transported supports at scale $k$, and let
$b>1$ be the block factor.  The theorem is deliberately formulated as a
certificate: an application must provide all constants and verify all
hypotheses without using clustering itself.

\begin{theorem}[Fibered two-tier clustering certificate, \conditional]
\label{thm:certificate}
Suppose a tower of $n$ fibered RG steps has the following properties for
centered local observables $F,G$ with initial support distance $d$.
\begin{enumerate}
\item \textbf{Quasilocal transport.}  The transported observables $F_k,G_k$
have norms at most $A_F,A_G$ and effective separation
$d_k\ge c_0d/b^k-c_1$.
\item \textbf{Good-fiber loss.}  For constants $C_f,m_f>0$ and $0\le K_k<1$,
\[
 \left|\E\Cov(F_k,G_k\mid\mathcal F_{k+1});\mathcal G_k\right|
 \le C_f A_FA_G K_k e^{-m_fd_k}.
\]
\item \textbf{Defect channel.}  Connected bad objects joining the two
transported supports contribute at most
\[
 C_bA_FA_G R_k e^{-m_bd_k},
\]
where $R_k\ge0$ is established by a dependent polymer or domination estimate.
\item \textbf{Terminal base rigidity.}  For some $C_*,m_*>0$,
\[
 |\Cov(F_n,G_n)|\le C_*A_FA_Ge^{-m_*d_n}.
\]
\end{enumerate}
Then
\begin{equation}\label{eq:certificate-bound}
 |\Cov(F,G)|\le A_FA_G\left[
 C_*e^{-m_*d_n}+
 \sum_{k=0}^{n-1}(C_fK_k e^{-m_fd_k}+C_bR_ke^{-m_bd_k})
 \right].
\end{equation}
In particular, any cutoff-uniform choice of $n=n(d)$ for which the right-hand
side is bounded by $Ce^{-m d}$ proves exponential clustering on the specified
observable core.
\end{theorem}
\begin{proof}
Take absolute values in \eqref{eq:iterated}.  Split each conditional fiber term
into its good and bad contributions and apply hypotheses 2 and 3.  Apply
hypothesis 4 to the terminal covariance and sum.  The last sentence is exactly
the stated uniform bound.
\end{proof}

The theorem does not say that $K_k+R_k<1$ alone automatically gives physical
exponential decay; scale conversion, support growth, amplitudes, and the number
of steps matter.  Its value is logical: retained information appears only in
the transported terminal term, while every eliminated degree of freedom appears
in a fiber increment.  No relevant mode is incorrectly charged as a mixing
failure.

\section{A restricted connected-defect theorem}

The simplest calibration makes bad blocks independent.  It cannot describe
four-dimensional gauge defects, but it gives a falsifiable numerical target.

\begin{theorem}[Independent rooted defect clusters, \proved]
\label{thm:defect}
Let $\Gamma$ be a graph of maximum degree $\Delta$, and mark vertices bad
independently with probability $p$.  For a root $o$ and $n\ge1$,
\[
 \Pr\{\text{there exists a connected all-bad }A\ni o, |A|=n\}
 \le \min\{1,\ p(\Delta^2p)^{n-1}\}.
\]
Thus $\Delta^2p<1$ gives exponential decay in $n$ for this event.
\end{theorem}
\begin{proof}
Every finite connected $A\ni o$ has a spanning tree.  A depth-first traversal
of that tree starting at $o$ has at most $2(n-1)$ directed edge steps.  There
are at most $\Delta^{2(n-1)}$ such walks, hence at most that many encoded
connected sets (overcounting is harmless).  A fixed set of $n$ distinct
vertices is all bad with probability $p^n$.  The union bound gives
$\Delta^{2(n-1)}p^n=p(\Delta^2p)^{n-1}$.
\end{proof}

In the Gaussian bridge of Theorem~\ref{thm:gaussian}, declaring the centered
fiber bad when $|R|>\rho\sqrt{t/2}$ gives
\[
 p_\rho=\Pr\{|N(0,1)|>\rho\}
 =\operatorname{erfc}(\rho/\sqrt2).
\]
At $\rho=3$, $p_3\approx0.0026998$.  With $\Delta=8$,
$\Delta^2p_3\approx0.1728<1$.  This proves subcriticality only for the
independent Gaussian calibration.  In Yang--Mills theory, bad blocks are
dependent and may contain large curvature, cut-locus effects, exceptional gauge
orbits, or topological flux.  A valid $R_k$ must control those mechanisms rather
than assume them away.

\section{Regulator comparison at block level}

\begin{theorem}[Finite-product TV comparison, \proved]
\label{thm:tensor}
For probability measures $\mu,\nu$ and integer $m\ge1$,
\[
 \TV(\mu^{\otimes m},\nu^{\otimes m})
 \le 1-(1-\TV(\mu,\nu))^m
 \le m\,\TV(\mu,\nu).
\]
\end{theorem}
\begin{proof}
Choose a maximal coupling of one $\mu$ and one $\nu$ sample, whose mismatch
probability is $\TV(\mu,\nu)$.  Couple $m$ pairs independently.  The product
samples agree unless at least one pair mismatches, which has probability
$1-(1-\TV)^m$.  Total variation is no larger than the mismatch probability of
any coupling, and Bernoulli's inequality gives the second bound.
\end{proof}

Consequently convergence of heat-kernel and Wilson one-plaquette laws implies
agreement for each fixed finite product block.  It does \emph{not} imply
universality for a fixed physical block containing $m(a)\to\infty$ plaquettes:
the displayed bound would require the one-plaquette discrepancy to be
$o(1/m(a))$, unless correlations or a sharper comparison provide compensation.
The same warning applies more strongly to interacting nonabelian blocks.

\section{Executable compact calibration}\label{sec:numerics}

The accompanying program evaluates the compact bridge on a periodic grid,
using the median-corrected cumulative-distribution formula for geodesic
$W_1$ on the circle.  Boundary increments are $0.1\sqrt t$; defects use
$\rho=3$.  The Wilson one-plaquette density is matched to the heat-kernel
quadratic small-angle scale before TV comparison.

\begin{center}
\small
\begin{tabular}{@{}rrrrrrrr@{}}
\toprule
$t$ & $\delta s$ & full $W_1/\delta s$ & fiber $W_1/\delta s$
& $p_{\rm bad}$ & $64p_{\rm bad}$ & one TV & $16\,$TV bound\\
\midrule
0.50 & 0.07071 & 0.498316 & $1.8\!\times\!10^{-8}$ & 0.002696 & 0.172523 & 0.063595 & 1.000000\\
0.20 & 0.04472 & 0.499999 & $<10^{-9}$ & 0.002697 & 0.172636 & 0.019918 & 0.318686\\
0.10 & 0.03162 & 0.500000 & $<10^{-9}$ & 0.002688 & 0.172036 & 0.009232 & 0.147705\\
0.05 & 0.02236 & 0.500000 & $<10^{-9}$ & 0.002719 & 0.174015 & 0.004477 & 0.071629\\
\bottomrule
\end{tabular}
\end{center}

Eight automated tests verify normalization, the $1/2$ tangent response, small
centered response, the Gaussian-tail limit, defect subcriticality in the stated
model, TV tensorization, and failure-closed input checks.  These computations
are reproducible evidence about the selected finite model, not computer-assisted
proofs for nonabelian four-dimensional Yang--Mills theory.

\section{The theorem dependency chain}

The representational programme now factors the Clay target as follows:
\begin{equation}\label{eq:chain}
\begin{split}
&\text{regulated gauge measure}
\xrightarrow{\text{gauge-covariant block/disintegration}}
\text{base + conditional fiber}\\
&\xrightarrow{\ K_k\ \text{and}\ R_k\ \text{bounds}\ }
\text{uniform fibered basin}
\xrightarrow{\text{Theorem \ref{thm:certificate}}}
\text{Euclidean clustering}\\
&\xrightarrow{\text{continuum coherence + OS axioms}}
\text{physical Hilbert space}
\xrightarrow{\text{spectral transfer}}
\operatorname{spec}(H)\subset\{0\}\cup[\Delta,\infty).
\end{split}
\end{equation}

Theorems \ref{thm:gaussian}--\ref{thm:tensor} clarify the middle language but
do not supply the four-dimensional arrows in \eqref{eq:chain}.  In particular,
the existence of a cutoff sequence and compactness alone are insufficient for
the stated endpoint if nontriviality, axioms, and a physical-unit gap are not
preserved \cite{jaffe-witten}.

\section{Frontier audit: what remains open}

\begin{longtable}{@{}p{0.08\textwidth}p{0.58\textwidth}p{0.22\textwidth}@{}}
\toprule
ID & Required result & Status\\
\midrule
\endhead
O1 & Construct local, gauge-covariant four-dimensional block maps and
disintegrations; specify retained loop data, fiber quotients, stabilizers,
overlap consistency, locality, and reflection compatibility. & \open\\
O2 & On good blocks prove a quotient Wasserstein or related influence bound
$K_k<1$, uniform in volume and cutoff, after retained modes are removed. & \open\\
O3 & Control dependent large-field, exceptional-orbit, cut-locus, and
topological defects by a connected polymer estimate with summable $R_k$. & \open\\
O4 & Prove that the asymptotically free ultraviolet family enters the fibered
basin with constants in physical units, without defining the basin by the
desired clustering conclusion. & \open\\
O5 & Prove convergence and regulator independence on a separating
gauge-invariant observable core, together with reflection positivity,
Euclidean covariance, symmetry, regularity, and nontriviality. & \open\\
O6 & Reconstruct the Hilbert space and show that uniform Euclidean decay gives
$\operatorname{spec}(H)\subset\{0\}\cup[\Delta,\infty)$ with $\Delta>0$. & \open\\
\bottomrule
\end{longtable}

The frontier theorem is O4, but it should not be attacked before O1--O3 are
made quantitative.  The immediate next nonabelian experiment is the smallest
$SU(2)$ block with explicitly described retained conjugacy/loop data.  It should
be tested against noncommuting boundary changes and four adversarial defect
classes: large curvature, holonomy near the cut locus, reducible stabilizers,
and nontrivial flux.  A positive numerical coefficient is useful only after the
metric, quotient, and locality constants are fixed independently of the target
claim.

\section{Conclusion}

The new ontology changes the proof obligation in a concrete way.  An RG step is
not a map that should erase all information.  It is a typed transformation that
\emph{preserves} selected relational information in its base and
\emph{forgets} conditional fiber information with controlled error.  The exact
covariance decomposition records both operations.  The mass-gap programme then
asks whether infrared rigidity is stable across this chain and survives the
continuum/reconstruction interface.

The two-plaquette bridge supplies a sharp sanity check: failure of full-mode
contraction can coexist with perfect fiber contraction.  This rescues a viable
local target but does not establish it for four-dimensional Yang--Mills.  The
remaining mathematical core is now narrower and more honest: construct the
nonabelian conditional fibers, prove their good-field contraction, control the
dependent defect channel, and carry the resulting physical-unit estimate
through a coherent continuum limit.

\begin{thebibliography}{9}
\bibitem{jaffe-witten}
A.~Jaffe and E.~Witten,
\emph{Quantum Yang--Mills Theory}, Clay Mathematics Institute Millennium
Problem description (2000),
\url{https://www.claymath.org/wp-content/uploads/2022/06/yangmills.pdf}.

\bibitem{osterwalder-schrader1}
K.~Osterwalder and R.~Schrader,
``Axioms for Euclidean Green's functions,''
\emph{Communications in Mathematical Physics} \textbf{31} (1973), 83--112.

\bibitem{osterwalder-schrader2}
K.~Osterwalder and R.~Schrader,
``Axioms for Euclidean Green's functions II,''
\emph{Communications in Mathematical Physics} \textbf{42} (1975), 281--305.

\bibitem{driver2d}
B.~K. Driver,
``Two dimensional Yang--Mills theory: expectations for Wilson loops,''
heat-kernel/Wilson-loop notes,
\url{https://www.math.ucsd.edu/~bdriver/DRIVER/Papers/Drivers_Papers/A3-YM2-Expectations.pdf}.
\end{thebibliography}

\end{document}
